\documentclass[10pt,a4paper]{article}
\usepackage[T1]{fontenc}
\usepackage[utf8]{inputenc}
\usepackage{lmodern}
\usepackage[italian]{babel}
\usepackage[margin=1.55cm]{geometry}
\usepackage{graphicx,booktabs,tabularx,enumitem,microtype,xcolor}
\usepackage[hidelinks]{hyperref}
\hypersetup{pdftitle={Architettura di Tandem MCP 4.2.0},pdfauthor={Alessandro Nocentini (ALENOC)},pdflang={it}}
\setlist{nosep,leftmargin=1.2em}
\setlength{\parindent}{0pt}
\setlength{\parskip}{2.8pt}
\newcommand{\code}[1]{\texttt{\detokenize{#1}}}

\title{\vspace{-1.8em}Architettura di Tandem MCP\\\large Versione 4.2.0: orchestratore cloud, esecutore locale}
\author{Alessandro Nocentini (ALENOC)\\\small\url{https://github.com/ALENOC/tandem-mcp}}
\date{Novembre 2026}

\begin{document}
\maketitle
\vspace{-1.2em}

\section{Panoramica}
Tandem MCP è un server Model Context Protocol che permette al modello della chat di affidare parte di un
lavoro di programmazione a un modello che gira sul proprio computer. L'orchestratore cloud (Claude Code,
Codex, Antigravity, Gemini CLI o un altro client MCP) tiene progetto, contratti e revisione. Il modello
locale (LM Studio, Ollama, llama.cpp, vLLM o un server compatibile OpenAI) legge i file grandi e scrive il
codice entro i limiti del contratto, e Tandem controlla ogni risultato prima che arrivi all'orchestratore.

\begin{center}
\includegraphics[width=0.52\textwidth]{diagram-it.pdf}
\end{center}

\section{Moduli principali}
\begin{tabularx}{\textwidth}{@{}lX@{}}
\toprule
Modulo & Compito\\
\midrule
\code{server.py} & 20 tool MCP (4 nel profilo autonomo), validazione rigorosa del contratto, ricevute, screenshot \code{ImageContent}, oscuramento segreti\\
\code{engine.py} & Ricerca del server, verifica delle capacità, scelta del modello per ruolo (digest, delegate, vision), retry con backoff\\
\code{digest.py} & Lettura a blocchi map-reduce, verifica lessicale delle citazioni, numeri non supportati, unione per schema\\
\code{agent.py} & Ciclo esecutore a prefisso stabile, compattazione lazy, \code{replace_in_file} multiplo, registro, lotti, rollback\\
\code{sandbox.py} & Glob, parser dei comandi, rifiuto dei symlink esterni, impronta SHA-256, jail Bubblewrap\\
\code{ui_test.py} & Scenari \code{tandem:ui_test} del contratto, eseguiti in Chromium headless\\
\code{web_check.py} & Controllo integrato di esecuzione, animazione, soglia 40 FPS e cattura screenshot 1280x800 per pagine web\\
\code{visual_review.py} & Gestione \code{reference_image}, critica visiva stateless con VLM locale (0 token cloud) e codifica screenshot per MCP\\
\code{router.py} & \code{auto_route}: strumenti deterministici, risposta diretta o delega locale\\
\code{host_runner.py} & Attesa del lavoro gestita da un processo locale, senza interrogare il cloud\\
\code{store.py} & Persistenza in \code{TANDEM_HOME}, lock tra processi, riconciliazione all'avvio\\
\bottomrule
\end{tabularx}

\section{Delega, revisione visiva e profilo autonomo}
\begin{itemize}
\item \textbf{Contratto:} \code{delegate} richiede un obiettivo, \code{write_globs}, \code{verify_commands} che
  falliscono finché il lavoro non è fatto e vincoli atomici, un requisito ciascuno (più un eventuale
  \code{reference_image} per lavori UI). Nel profilo autonomo il server divide i vincoli composti invece di rifiutare.
\item \textbf{Controllo iniziale e protezione:} i controlli girano prima di ogni modifica; i file di test sono in sola
  lettura e vengono ripristinati se toccati.
\item \textbf{Ciclo locale e riparazione:} la conversazione cresce solo in coda così il server locale riusa la
  cache del prompt; ogni risposta arriva a 12.000 token (\code{TANDEM_REPLY_TOKENS}). A \code{finish} Tandem
  controlla sintassi, controlli da shell nella sandbox (rimandando una volta quelli lenti oltre
  \code{TANDEM_SLOW_CHECK_S}), \code{tandem:ui_test}, soglia 40 FPS in Chromium headless (salvando uno screenshot
  JPEG 1280x800 nella cartella del lavoro) e righe citate. L'output tagliato conserva le righe fallite; tre
  fallimenti identici di fila fermano il lavoro come \code{stuck}.
\item \textbf{Revisione visiva (0 codice cloud):} superati i controlli funzionali e gli FPS, se è disponibile un VLM
  locale (\code{TANDEM_MODEL_VISION} o un VLM già caricato con \code{TANDEM_VISUAL_REVIEW=auto}), \code{visual_review.py}
  esegue una critica stateless confrontando lo screenshot con contratto ed eventuale \code{reference_image}, rimandando
  correzioni concrete nel ciclo locale (\code{TANDEM_MAX_VISUAL_REPAIRS=2}) senza alterare la cache KV del coding. Con
  \code{TANDEM_RECEIPT_SCREENSHOT=1} le ricevute web allegano lo screenshot come \code{ImageContent} MCP: l'orchestratore
  può verificare la resa e mandare un breve \code{job_continue(hint)} senza scrivere codice.
\item \textbf{Profilo autonomo} (\code{TANDEM_PROFILE=autonomous}): quattro tool. \code{delegate} resta in
  attesa in locale senza chiamate cloud, entro \code{TANDEM_JOB_SECONDS} (3600), \code{TANDEM_MAX_REPAIRS} (20)
  e \code{TANDEM_MAX_STEPS} (160). Notifiche di avanzamento ogni 30 s evitano timeout del client.
\end{itemize}

\section{Sicurezza e rollback}
\begin{itemize}
\item \textbf{Profili della sandbox} (\code{TANDEM_SANDBOX}): \code{secure} (si blocca se manca l'isolamento)
  di default su Linux, \code{compatible} su macOS e Windows, \code{trusted} disattiva Bubblewrap.
\item \textbf{Jail Bubblewrap:} cartelle di sistema e runtime Python sono montati in sola lettura e il workspace
  è l'unica cartella scrivibile; \code{$HOME}, \code{/tmp} e \code{/run} dell'host sono vuoti; PID e IPC sono
  separati, e anche la rete salvo che il compito la consenta.
\item \textbf{Parser dei comandi:} ogni parte di una pipeline deve corrispondere ad \code{allow_run} e non
  alla lista dei divieti; sostituzione di comandi e heredoc sono rifiutati.
\item \textbf{Backup e rollback:} i byte originali sono registrati per ogni lavoro (8 MiB per file, 256 MiB
  per lavoro); un rollback incompleto imposta lo stato \code{partial_rollback}.
\item \textbf{Oscuramento dei segreti:} l'output dei tool nasconde chiavi private, token e variabili segrete.
\end{itemize}

\section{Risparmio misurato}
Con Opus 5.5 come orchestratore e il profilo autonomo (campagna 4.2.0 \code{v4_2}, cinque compiti, tre coppie
bilanciate ciascuno, 15 tentativi per braccio), entrambi i bracci hanno ottenuto 10/15 risultati nel budget e Tandem
ha ridotto la spesa cloud totale e il costo per successo del 53\% (\$0,202 contro \$0,429), con il 68\% di token di
output in meno (-49\% per successo escludendo uno stallo API cloud di 82 minuti sul braccio diretto). I compiti medi
e grandi hanno risparmiato dal 22\% al 51\% per successo e il portale 3D è costato il 68\% in meno per portale consegnato; la
correzione di 5 righe costa l'8\% in più e resta diretta. Il tempo mediano di Tandem sui compiti comuni è sceso da
25 min (4.0.0) a 19 min (4.1.0) a 16 min, ancora circa 12 volte più lento del diretto. Tutte le pagine Tandem
promosse girano a 60 FPS. Il VLM locale ha revisionato 7 pagine (2 segnalazioni, entrambe falsi positivi) e Opus ha
inviato 3 suggerimenti visivi dopo aver visto lo screenshot della ricevuta. Modello locale: Qwen3.8 27B IQ3\_S con
\code{mmproj-F16} (\code{type=vlm}, contesto 40.960 token) su GPU AMD Radeon RX 6800 serie (16 GB). Dettagli:
\code{docs/BENCHMARKS.md} e \url{https://alenoc.github.io/tandem-mcp/}.

\end{document}
